\section{Introduction}
\label{sec:intro}

A large-language-model serving deployment is, from the network's point of view,
a small number of very regular traffic patterns repeated millions of times:
tensor-parallel all-reduces between a handful of ranks, expert dispatch and
combine in a mixture-of-experts layer, key-value cache movement between a
prefill instance and a decode instance. Each of them is a short burst on one or
a few queue pairs, separated by a compute gap of a few microseconds to a few
hundred microseconds. Whether such a burst runs at line rate, at half of it, or
at a rate set by a retransmission equilibrium is decided inside the RDMA NIC and
by its interaction with the fabric. A simulator that predicts time to first
token and time per output token therefore inherits the NIC's behaviour whether
it models it or not.

The obstacle is that no public contract for that behaviour exists. The
InfiniBand Architecture Specification and its RoCEv2 annex define wire formats
and transport semantics but say nothing about how many work queue entries a
device keeps in flight, how deep its ingress buffer is, or when it decides to
drop. The vendor's programmer's reference manual documents a software
interface, not a microarchitecture. The open-source \texttt{mlx5} driver and
the \texttt{rdma-core} provider show exactly how work is submitted and
completions are read, which is a great deal, and nothing at all about what the
device does in between. Published RNIC designs are either research prototypes
with different design points or commercial IP with an NDA attached. The result
is that everyone who needs an endpoint model either fits one number (a fixed
per-message latency) and hopes, or reuses somebody else's constant from a
different generation of silicon.

We take the other route. We treat the device as a black box with a
well-documented interface, we probe it under conditions chosen so that a
mechanism becomes identifiable rather than merely visible, and we turn the
result into an architecture that can be built. The target of the build is an
AMD Alveo U55C accelerator card, which carries the two 100\,GbE ports, the PCIe
attachment and the high-bandwidth memory that an endpoint of this class needs;
the architecture is parameterised so that the same register-transfer-level
sources move to an ASIC by exchanging the memory and hard-macro choices rather
than the datapath.

\subsection{Why a measurement-calibrated, open, implementable model}

Three consumers want the same artifact for three different reasons.

\textbf{Simulator fidelity.} SimLLM, the serving simulator this work feeds,
predicts the behaviour of continuous-batching serving
schedulers~\cite{vllm2023}, and therefore needs an endpoint whose timing is
right at the granularity of a single collective. A model that reproduces line
rate for a large message and nothing else will predict a tensor-parallel decode step wrongly by a factor that
changes with queue depth, message size and whether the previous burst left the
receiver's ingress buffer drained. Our measurements show all three effects and
quantify them.

\textbf{A golden reference for RTL.} An open RDMA endpoint implementation needs
something to be verified against. A specification is not enough, because the
specification does not fix the timing or the counter semantics; a second
implementation is circular. A deterministic C model, calibrated against
measured silicon and carrying an explicit table of the behaviours it is
required to reproduce, is a usable oracle: the same stimulus goes through the
model and through the design under test, and the comparison is over timestamps
and counters rather than over a bandwidth curve that looks about right.

\textbf{A public account of what the silicon does.} Several of the behaviours
below are not in any datasheet and would surprise a reader who has only read
the specification. The NIC overwrites the two explicit-congestion-notification
bits of every RoCEv2 packet it sends. The standard ``is ECN happening'' counter
reads zero while congestion control is actively running. The counter that
appears to count lost packets counts loss bursts instead, at a ratio of about
73 to 1. A per-queue-pair receive cap that had been attributed to the NIC for
most of the campaign turned out to belong to the measurement tool. Recording
these publicly, with the evidence and with the corrections that the campaign
made to its own earlier conclusions, is a contribution in its own right.

\subsection{Contributions}

\begin{itemize}
  \item \textbf{A campaign method for edge-case probing on a lossy fabric}
    (\S\ref{sec:measurement}). Six phases, each with a frozen expectations
    file written before the run, more than 700 measured points, and an explicit
    validity review that marks which tests the fabric can and cannot support.
    Every parameter carries one of five evidence classes: \evdoc{}, \evdrv{},
    \evcal{}, \evdecl{} or \evinf{}.
  \item \textbf{The reverse-engineered mechanisms and their constants}
    (\S\ref{sec:measurement}). The fixed-offset initiation law and its two
    parameters, the queue-depth dominance, the loss equilibrium and the drain
    window, the internal processing budget and its arbitration order, the
    unconditional ECT(0) stamp, endpoint-generated congestion notification on a
    fabric that never marks, the tail-drop port buffer, and the counter
    semantics that any detection tool built on this hardware has to know.
  \item \textbf{A nineteen-block architecture with an AXI4 monitoring network
    on chip} (\S\ref{sec:arch}), targeting the Alveo U55C and expandable to an
    ASIC. Two interconnects: an AXI4 data path for host-memory and
    high-bandwidth-memory traffic, and a separate AXI4-Lite monitoring ring
    that gives every block a slot in one versioned register map mirroring the
    real NIC's observable state. Each block names the measured behaviour it
    must reproduce and the evidence class of every parameter it carries.
  \item \textbf{The golden C model and the anomaly table as a verification
    contract} (\S\ref{sec:verify}). Nineteen registered anomalies, each with a
    kind (\textit{emergent}, \textit{injected}, \textit{fabric},
    \textit{counter} or \textit{tool}), a trigger, a magnitude and a band. A
    block is validated when the anomalies registered against it are reproduced
    inside their bands, and not before.
\end{itemize}

\subsection{Scope and honesty about it}

Everything here concerns one device on one fabric: ConnectX-5 Ex at 100\,GbE,
RoCEv2, port maximum transmission unit 4096, on a cluster with priority flow
control disabled, global pause enabled but absorbed one hop away, and DCQCN
running at the firmware default. Some findings are properties of the silicon
(the initiation law, the internal budget, the ECT(0) stamp); some are
properties of the deployment (the tail-drop buffer, the absence of switch
marking); and the paper marks which is which, because a model that confuses
them will mispredict on any other cluster. Where a later phase corrected an
earlier conclusion we state the correction and keep the earlier claim visible,
because the correction history is the method working, not a defect in it.

We also make no claim to have recovered the vendor's design. What we recover is
an architecture that is consistent with everything we could observe from
outside, that reproduces the measurements within stated bands, and that can be
built. Several blocks contain a fitted stand-in where the real mechanism is not
identifiable from the outside, and each of those is labelled \evcal{} rather than
dressed up as a register.
